\documentclass[11pt,a4paper]{article}

% =============================================================================
% FEHRADVICE STYLE - Based on Stadionumbau Kommunikationskonzept
% Clean, minimal design with blue accents
% =============================================================================

% Packages
\usepackage[utf8]{inputenc}
\usepackage[T1]{fontenc}
\usepackage[german]{babel}
\usepackage[
    top=25mm,
    bottom=25mm,
    left=25mm,
    right=25mm
]{geometry}
\usepackage{graphicx}
\usepackage{xcolor}
\usepackage{colortbl}
\usepackage{booktabs}
\usepackage{tabularx}
\usepackage{array}
\usepackage{fancyhdr}
\usepackage{lastpage}
\usepackage{enumitem}
\usepackage{hyperref}
\usepackage{helvet}
\usepackage{tikz}
\usepackage{amssymb}
\usepackage{titlesec}
\usepackage{parskip}
\renewcommand{\familydefault}{\sfdefault}

% Fix headheight
\setlength{\headheight}{26pt}
\addtolength{\topmargin}{-6pt}
\setlength{\parindent}{0pt}
\setlength{\parskip}{10pt}

% =============================================================================
% COLORS
% =============================================================================

\definecolor{fadarkblue}{RGB}{2,64,121}
\definecolor{falightblue}{RGB}{84,158,222}
\definecolor{fadarkgray}{RGB}{37,33,42}
\definecolor{falightgray}{RGB}{243,245,247}
\definecolor{spoered}{RGB}{206,0,0}

% =============================================================================
% HYPERREF
% =============================================================================

\hypersetup{
    colorlinks=true,
    linkcolor=fadarkblue,
    citecolor=fadarkblue,
    urlcolor=fadarkblue
}

% =============================================================================
% HEADER & FOOTER - FehrAdvice Style
% =============================================================================

\pagestyle{fancy}
\fancyhf{}
\fancyhead[L]{\small\color{fadarkgray}SEITE \thepage\ VON \pageref{LastPage} — \textbf{SPÖ PROJEKT 2029} — FEHRADVICE \& PARTNERS AG}
\fancyhead[R]{\begin{tikzpicture}[scale=0.3]
    \fill[fadarkblue] (0,0) rectangle (0.3,1);
    \fill[fadarkblue] (0.4,0) rectangle (0.7,1);
    \fill[fadarkblue] (0.8,0) rectangle (1.1,1);
    \fill[fadarkblue] (1.2,0) rectangle (1.5,1);
\end{tikzpicture}}
\renewcommand{\headrulewidth}{0pt}
\renewcommand{\footrulewidth}{0pt}

% First page style (no header)
\fancypagestyle{plain}{
    \fancyhf{}
    \fancyhead[L]{\small\color{fadarkgray}SEITE \thepage\ VON \pageref{LastPage} — \textbf{SPÖ PROJEKT 2029} — FEHRADVICE \& PARTNERS AG}
    \fancyhead[R]{\begin{tikzpicture}[scale=0.3]
        \fill[fadarkblue] (0,0) rectangle (0.3,1);
        \fill[fadarkblue] (0.4,0) rectangle (0.7,1);
        \fill[fadarkblue] (0.8,0) rectangle (1.1,1);
        \fill[fadarkblue] (1.2,0) rectangle (1.5,1);
    \end{tikzpicture}}
    \renewcommand{\headrulewidth}{0pt}
}

% Title page style
\fancypagestyle{titlepage}{
    \fancyhf{}
    \fancyhead[L]{\small\color{fadarkgray}SEITE \thepage\ VON \pageref{LastPage} — \textbf{SPÖ PROJEKT 2029} — FEHRADVICE \& PARTNERS AG}
    \fancyhead[R]{\begin{tikzpicture}[scale=0.4]
        \fill[fadarkblue] (0,0) rectangle (0.3,1);
        \fill[fadarkblue] (0.4,0) rectangle (0.7,1);
        \fill[fadarkblue] (0.8,0) rectangle (1.1,1);
        \fill[fadarkblue] (1.2,0) rectangle (1.5,1);
        \node[right] at (1.7,0.5) {\color{fadarkblue}\sffamily\textbf{FEHR}};
        \node[right] at (1.7,0) {\color{fadarkblue}\sffamily\textbf{ADVICE}};
    \end{tikzpicture}}
    \renewcommand{\headrulewidth}{0pt}
}

% =============================================================================
% SECTION FORMATTING - Large number left, title right
% =============================================================================

\titleformat{\section}
    {\sffamily\bfseries\Large\color{fadarkblue}}
    {\thesection}
    {1em}
    {\MakeUppercase}

\titleformat{\subsection}
    {\sffamily\bfseries\normalsize\color{fadarkgray}}
    {}
    {0em}
    {}

\titlespacing*{\section}{0pt}{20pt}{10pt}
\titlespacing*{\subsection}{0pt}{15pt}{5pt}

% =============================================================================
% LIST STYLING - Simple black bullets
% =============================================================================

\setlist[itemize]{
    label=\textbullet,
    leftmargin=2em,
    itemsep=2pt,
    parsep=0pt
}

\setlist[enumerate]{
    leftmargin=2em,
    itemsep=2pt,
    parsep=0pt
}

% =============================================================================
% TABLE STYLING - Simple with blue header
% =============================================================================

\newcolumntype{L}[1]{>{\raggedright\arraybackslash}p{#1}}
\newcolumntype{C}[1]{>{\centering\arraybackslash}p{#1}}

% =============================================================================
% DOCUMENT
% =============================================================================

\begin{document}

\color{fadarkgray}

% =============================================================================
% TITLE PAGE
% =============================================================================

\thispagestyle{titlepage}

\vspace*{2cm}

% Blue block with title
\noindent\begin{tikzpicture}
\fill[fadarkblue] (0,0) rectangle (\textwidth,-12cm);
\node[anchor=north west, text=white, font=\sffamily\bfseries\Huge] at (1cm,-1.5cm) {ORDNEN};
\node[anchor=north west, text=white, font=\sffamily\bfseries\Huge] at (1cm,-3cm) {STATT};
\node[anchor=north west, text=white, font=\sffamily\bfseries\Huge] at (1cm,-4.5cm) {SPALTEN};
\node[anchor=north west, text=white, font=\sffamily\large] at (1cm,-6.5cm) {Strategisches Kommunikationsbriefing};
\node[anchor=north west, text=white, font=\sffamily\large] at (1cm,-7.3cm) {SPÖ Projekt 2029};
\node[anchor=north west, text=white, font=\sffamily] at (1cm,-9cm) {FehrAdvice \& Partners AG};
\node[anchor=north west, text=white, font=\sffamily] at (1cm,-9.7cm) {Jänner 2026};
\end{tikzpicture}

\vfill

\begin{center}
{\Large\color{spoered}\textbf{VERTRAULICH}}

\vspace{0.3cm}

{\small Nur für benannte Empfänger:innen}
\end{center}

\newpage

% =============================================================================
% INHALT (Table of Contents)
% =============================================================================

\section*{INHALT}

\begin{tabular}{@{}l @{\hspace{1cm}} r@{}}
1 & Executive Summary \dotfill\ 3 \\[8pt]
2 & Sachverhalt und Fragestellung \dotfill\ 4 \\[8pt]
3 & Verhaltensökonomische Analyse \dotfill\ 5 \\[8pt]
4 & Die drei strategischen Optionen \dotfill\ 6 \\[8pt]
5 & Operative Umsetzung: 10 Kontexte \dotfill\ 8 \\[8pt]
6 & Die Parteitagsrede \dotfill\ 9 \\[8pt]
7 & Übersichtstabellen \dotfill\ 10 \\[8pt]
APPENDIX & Leitmotiv \& Checklisten \dotfill\ 11 \\
\end{tabular}

\newpage

% =============================================================================
% 1. EXECUTIVE SUMMARY
% =============================================================================

\section{EXECUTIVE SUMMARY}

\textbf{Die SPÖ steht vor der Wahl: Jeden Angriff abwehren (und verlieren) oder das Spielfeld ändern (und gewinnen).}

«Ordnen statt Spalten» ist nicht nur ein Slogan -- es ist eine komplette Repositionierung der österreichischen Sozialdemokratie.

\subsection{Ausgangslage}

Am 25. Jänner 2026 veröffentlichte die Kronen Zeitung einen Artikel mit dem Titel «Spitals-Touristen kosten das Gesundheitssystem Milliarden». Der Artikel bedient ein von der FPÖ getriebenes Narrativ.

\textbf{Der Killer-Fakt:}

\begin{center}
\fbox{\parbox{0.7\textwidth}{\centering\large\textbf{2,7\% der Behandlungen} bei \textbf{4,8\% Bevölkerungsanteil}\\[5pt]= UNTERPROPORTIONALE NUTZUNG}}
\end{center}

\subsection{Empfehlung: Phasenweise Eskalation}

\begin{tabular}{@{}>{\columncolor{fadarkblue}\color{white}\bfseries}c L{2.5cm} L{4.5cm} L{2.5cm}@{}}
\rowcolor{fadarkblue}
\textcolor{white}{\textbf{Phase}} & \textcolor{white}{\textbf{Zeitraum}} & \textcolor{white}{\textbf{Primärstrategie}} & \textcolor{white}{\textbf{Intensität}} \\
\midrule
1 & 27. Jan -- 15. Feb & Ordnungs-Pivot & Defensiv-neutral \\
2 & 16. Feb -- 6. März & Ordnungs-Pivot + Spalter-Entlarvung & Aufbauend \\
3 & 7. März & Volle Offensive & Parteitagsrede \\
\end{tabular}

\subsection{Das Leitmotiv: Populär statt populistisch}

\begin{tabular}{@{}>{\columncolor{fadarkblue}\color{white}\bfseries}l L{4cm} L{4.5cm}@{}}
\rowcolor{fadarkblue}
\textcolor{white}{\textbf{Dimension}} & \textcolor{white}{\textbf{Populistisch (FPÖ)}} & \textcolor{white}{\textbf{Populär (SPÖ neu)}} \\
\midrule
Emotion & Angst, Wut, Hass & Sicherheit, Stolz, Zugehörigkeit \\
Einfachheit & Lügen vereinfachen & Wahrheit vereinfachen \\
Feind & Minderheiten & Spaltung selbst \\
Lösung & Ausgrenzung & Ordnung \\
Langfristeffekt & Zerstört Gesellschaft & Stärkt Gesellschaft \\
\end{tabular}

\vspace{10pt}

\textit{«Ordnen statt Spalten» ist populär, nicht populistisch. Es ist emotional, ohne zu spalten. Es ist einfach, ohne zu lügen. Es ist kraftvoll, ohne zu zerstören.}

\newpage

% =============================================================================
% 2. SACHVERHALT
% =============================================================================

\section{SACHVERHALT UND FRAGESTELLUNG}

\subsection{Der Anlass: Krone-Artikel vom 25. Jänner 2026}

\textbf{Titel:} «Spitals-Touristen kosten das Gesundheitssystem Milliarden»

\begin{tabular}{@{}>{\columncolor{fadarkblue}\color{white}}L{3cm} L{7cm}@{}}
\rowcolor{fadarkblue}
\textcolor{white}{\textbf{Behauptung}} & \textcolor{white}{\textbf{Fakt}} \\
\midrule
«22 Mio. Behandlungen» & Korrekt, aber über 10 Jahre und inkl. aller Kontakte \\
«Milliarden-Kosten» & Nicht belegt; 22 Mio. = 2,75\% von 800 Mio. Gesamtleistungen \\
Bevölkerungsanteil & 4,8\% -- d.h. \textbf{unterproportionale} Nutzung \\
«Schönheits-OPs» & Ästhetische Eingriffe werden NICHT öffentlich finanziert \\
\end{tabular}

\subsection{Der strategische Kontext: FPÖ-Boulevard-Achse}

\textbf{Das FPÖ-Playbook:}
\begin{enumerate}
    \item \textbf{Themen-Injection:} FPÖ-nahe Quellen liefern «Geschichten» an Boulevard
    \item \textbf{Emotionalisierung:} Grosse Zahlen ohne Kontext («22 Millionen!», «Milliarden!»)
    \item \textbf{Schuldattribution:} Implizite Verknüpfung: Migration $\rightarrow$ Kosten $\rightarrow$ Ihre Probleme
    \item \textbf{SPÖ-Falle:} SPÖ wird zur Reaktion gezwungen, verteidigt sich, verstärkt das Thema
\end{enumerate}

\subsection{Die zentrale Fragestellung}

\begin{center}
\fbox{\parbox{0.85\textwidth}{\centering\textbf{Wie verwandelt die SPÖ einen Boulevard-Angriff in eine Bestätigung ihrer eigenen strategischen Positionierung?}}}
\end{center}

\begin{tabular}{@{}>{\columncolor{fadarkblue}\color{white}}L{4.5cm} L{6cm}@{}}
\rowcolor{fadarkblue}
\textcolor{white}{\textbf{FPÖ-Frame}} & \textcolor{white}{\textbf{SPÖ-Gegen-Frame}} \\
\midrule
«Ausländer kosten euch Geld» & «Der Spaltungsboulevard kostet euch Vertrauen» \\
«SPÖ verteidigt Ausländer» & «SPÖ ordnet das Chaos der Vorgänger» \\
«Euer Steuergeld verschwindet» & «Euer Steuergeld finanziert nicht mehr den Spaltungsboulevard» \\
\end{tabular}

\newpage

% =============================================================================
% 3. VERHALTENSÖKONOMISCHE ANALYSE
% =============================================================================

\section{VERHALTENSÖKONOMISCHE ANALYSE}

\subsection{Das Grundproblem: Frame-Kontrolle}

\textit{«Wer den Frame kontrolliert, gewinnt die Debatte.»}

Die FPÖ operiert im \textbf{Ausschluss-Frame}:
\begin{itemize}
    \item «Wer gehört dazu?»
    \item «Wer kostet uns Geld?»
    \item «Wer ist schuld?»
\end{itemize}

Die SPÖ muss im \textbf{Funktions-Frame} bleiben:
\begin{itemize}
    \item «Was funktioniert?»
    \item «Was lösen wir?»
    \item «Wer bringt Ordnung?»
\end{itemize}

\subsection{Das Problem mit Verteidigung}

Wenn die SPÖ erklärt: «Die 22 Millionen sind eigentlich...»

$\rightarrow$ Impliziert: Die Frage ist berechtigt.\\
$\rightarrow$ Impliziert: Wir müssen uns rechtfertigen.\\
$\rightarrow$ Impliziert: Vielleicht stimmt ja doch was dran.

\textbf{Psychologisches Prinzip:} «Wer sich verteidigt, klagt sich an.»

\subsection{Kognitive Mechanismen}

\begin{tabular}{@{}>{\columncolor{fadarkblue}\color{white}}L{3cm} L{7cm}@{}}
\rowcolor{fadarkblue}
\textcolor{white}{\textbf{Mechanismus}} & \textcolor{white}{\textbf{Anwendung}} \\
\midrule
Verfügbarkeitsheuristik & Was zuletzt gesagt wird, bleibt im Gedächtnis. Bei Pivot zur Lösung («Kassenstellen») wird die Lösung erinnert. \\
Verlustaversion ($\lambda \approx 2.0$) & FPÖ aktiviert Verlust-Frame. SPÖ muss zum Gewinn pivoten ODER Verlust umlenken. \\
Attributionsfehler & FPÖ will: «Babler verteidigt Ausländer» (Person). SPÖ braucht: «Boulevard spaltet» (Situation). \\
\end{tabular}

\newpage

% =============================================================================
% 4. DIE DREI STRATEGIEN
% =============================================================================

\section{DIE DREI STRATEGISCHEN OPTIONEN}

\subsection{Strategie 1: Ordnungs-Pivot}

\textbf{Prinzip:} Minimale Faktennennung, maximaler Pivot zur eigenen Lösung.

\textbf{Formel:} Fakt (minimal) $\rightarrow$ Pivot $\rightarrow$ Lösung $\rightarrow$ Ende

\textbf{Beispiel-Wording:}

\begin{quote}
\textit{«2,7 Prozent der Behandlungen, bei 4,8 Prozent Bevölkerungsanteil. Das ist unter dem Durchschnitt. Das echte Problem sind fehlende Kassenstellen. Daran arbeiten wir.»}
\end{quote}

\begin{tabular}{@{}L{4cm} L{6cm}@{}}
\textbf{Risiko} & Niedrig \\
\textbf{Frame-Kontrolle} & Teilweise (50/50) \\
\textbf{Mobilisierungskraft} & Mittel \\
\textbf{Voraussetzungen} & Keine \\
\end{tabular}

\textbf{Wann einsetzen:} ZIB2, Ö3, breites Publikum. \textit{Default vor dem Parteitag.}

\subsection{Strategie 2: Spalter-Entlarvung}

\textbf{Prinzip:} Der Angriff selbst wird zum Beweis für die eigene These.

\textbf{Formel:} Angriff $\rightarrow$ «Genau das meinen wir» $\rightarrow$ Bestätigt Positionierung $\rightarrow$ Offensive

\textbf{Beispiel-Wording:}

\begin{quote}
\textit{«Genau das meinen wir mit Spaltungsboulevard. 2,7 Prozent werden zu ‹Milliarden› gemacht. Das ist keine Berichterstattung, das ist Spaltung. Und das Steuergeld der Österreicherinnen und Österreicher finanziert das nicht mehr. Das haben wir geordnet.»}
\end{quote}

\textbf{Die drei Frame-Elemente:}
\begin{enumerate}
    \item «Euer Steuergeld...» $\rightarrow$ Es war STEUERGELD, nicht Parteigeld
    \item «...nicht MEHR» $\rightarrow$ Die VORHERIGEN haben das gemacht
    \item «Das haben wir geordnet.» $\rightarrow$ SPÖ = Ordnung
\end{enumerate}

\textbf{Wann einsetzen:} Parteibasis, Gewerkschaft. \textit{Parteitagsrede: Volle Entfaltung.}

\subsection{Strategie 3: Agenda-Disruption}

\textbf{Prinzip:} Das gegnerische Thema wird durch ein eigenes, stärkeres Thema ersetzt.

\begin{quote}
\textit{«Ich rede heute über die 500 neuen Kassenstellen, die wir nächste Woche präsentieren werden. DAS ist Ordnung im Gesundheitssystem. Wer über Phantomzahlen reden will, kann das tun -- wir arbeiten.»}
\end{quote}

\textbf{Voraussetzung:} Nur mit Policy-Substanz! Ohne konkrete Ankündigung wirkt es wie Ausweichen.

\newpage

% =============================================================================
% 5. OPERATIVE UMSETZUNG
% =============================================================================

\section{OPERATIVE UMSETZUNG: 10 KONTEXTE}

\subsection{Das Babler-Spektrum}

Andreas Babler kann sowohl kämpferisch als auch staatsmännisch auftreten:

\begin{center}
\begin{tikzpicture}
  \draw[thick, <->] (0,0) -- (12,0);
  \node[below] at (0,-0.3) {\footnotesize KÄMPFERISCH};
  \node[below] at (12,-0.3) {\footnotesize STAATSMÄNNISCH};
  \node[above] at (1,0.2) {\scriptsize Parteitag};
  \node[above] at (3.5,0.2) {\scriptsize Social Media};
  \node[above] at (6,0.2) {\scriptsize ZIB2};
  \node[above] at (9,0.2) {\scriptsize Parlament};
  \node[above] at (11,0.2) {\scriptsize Opernball};
  \draw[fadarkblue, very thick] (6,-0.4) -- (6,0.6);
  \node[above, fadarkblue, font=\small\bfseries] at (6,0.7) {IMMER BABLER};
\end{tikzpicture}
\end{center}

\subsection{Kontext-Matrix}

\begin{tabular}{@{}>{\columncolor{fadarkblue}\color{white}\bfseries}c L{2.5cm} L{2.5cm} L{4cm}@{}}
\rowcolor{fadarkblue}
\textcolor{white}{\textbf{\#}} & \textcolor{white}{\textbf{Kontext}} & \textcolor{white}{\textbf{Strategie}} & \textcolor{white}{\textbf{Stil}} \\
\midrule
1 & Wirtshaus & S1/S2 & Dialekt, warm \\
2 & ZIB2 & S1 & Staatsmännisch, Fakten \\
3 & Ö3-Frühstück & S1 & Locker, ein Fakt, Pivot \\
4 & Twitter/X & S2 & Scharf, kurz, teilbar \\
5 & Parteitag & S2 voll & Kämpferisch, rhetorisch \\
6 & Parlament & S1/S3 & Formell, Akten-Pivot \\
7 & Pressekonferenz & S1/S3 & Sachlich, Policy-Substanz \\
8 & Gewerkschaft & S2 & Klassenkampf-Frame \\
9 & Dorffest & S1 & Persönlich, Lösungsorientiert \\
10 & Opernball & S1 minimal & Staatstragend, kein Angriff \\
\end{tabular}

\newpage

% =============================================================================
% 6. PARTEITAGSREDE
% =============================================================================

\section{DIE PARTEITAGSREDE}

\subsection{Crescendo-Struktur}

\begin{tabular}{@{}>{\columncolor{fadarkblue}\color{white}\bfseries}c L{2.5cm} L{7cm}@{}}
\rowcolor{fadarkblue}
\textcolor{white}{\textbf{Phase}} & \textcolor{white}{\textbf{Element}} & \textcolor{white}{\textbf{Inhalt}} \\
\midrule
1 & Einstieg & Persönliche Anekdote, emotionaler Anker \\
2 & Problem & Spaltung benennen, konkretes Beispiel \\
3 & Diagnose & «Ordnen statt Spalten» als Lösung \\
4 & Beweis & Steuergeld-Argument, Fakten \\
5 & Vision & Positives Österreich-Bild \\
6 & Appell & Call to Action, Mobilisierung \\
\end{tabular}

\subsection{Der rhetorische Höhepunkt}

\begin{quote}
\textit{«Sie haben dieses Land gespalten. Wir werden es ordnen.\\
Sie haben Angst geschürt. Wir werden Sicherheit schaffen.\\
Sie haben mit eurem Steuergeld den Hass finanziert. Wir haben das beendet.\\
Das ist der Unterschied zwischen Spalten und Ordnen.\\
Das ist der Unterschied zwischen ihnen und uns.\\
Das ist die SPÖ.»}
\end{quote}

\newpage

% =============================================================================
% 7. ÜBERSICHTSTABELLEN
% =============================================================================

\section{ÜBERSICHTSTABELLEN}

\subsection{Zielgruppensegmente}

\begin{tabular}{@{}>{\columncolor{fadarkblue}\color{white}}L{3cm} c L{2cm} L{4cm}@{}}
\rowcolor{fadarkblue}
\textcolor{white}{\textbf{Segment}} & \textcolor{white}{\textbf{Anteil}} & \textcolor{white}{\textbf{Hauptmotiv}} & \textcolor{white}{\textbf{Kommunikationscode}} \\
\midrule
SPÖ-Kernwähler & 20--25\% & Identität & «Wir sind die Guten» \\
Swing-Wähler Mitte & 25--30\% & Stabilität & «Ordnung bringt Ruhe» \\
Junge Wähler & 15\% & Zukunft & «Keine Angst, Lösungen» \\
Ältere Wähler & 20\% & Sicherheit & «Wir schützen, was zählt» \\
FPÖ-Skeptische & 10--15\% & Enttäuschung & «Es gibt einen anderen Weg» \\
\end{tabular}

\subsection{Risiken und Antworten}

\begin{tabular}{@{}>{\columncolor{fadarkblue}\color{white}}L{4cm} L{6cm}@{}}
\rowcolor{fadarkblue}
\textcolor{white}{\textbf{Risiko / Kritik}} & \textcolor{white}{\textbf{Antwort}} \\
\midrule
«Babler ist zu links» & «Ich bin für Ordnung. Das ist nicht links oder rechts, das ist vernünftig.» \\
«SPÖ ist migrantenfreundlich» & «Ich bin für faire Regeln. Für alle. Das ist Ordnung.» \\
«Wo sind die konkreten Lösungen?» & «500 neue Kassenstellen. Das ist konkret.» \\
«Warum greift ihr die Krone an?» & «Wir greifen niemanden an. Wir haben das Steuergeld umgeleitet.» \\
\end{tabular}

\newpage

% =============================================================================
% APPENDIX
% =============================================================================

\section*{APPENDIX: LEITMOTIV \& CHECKLISTEN}

\subsection*{Das Leitmotiv erklärt}

\textbf{«Populär statt populistisch»}

Die SPÖ kann genauso emotional und eingängig kommunizieren wie die FPÖ -- ohne zu lügen, zu spalten oder zu zerstören. Das ist der dritte Weg.

\begin{tabular}{@{}L{2.5cm} L{4cm} L{4cm}@{}}
\textbf{Aspekt} & \textbf{Populistisch} & \textbf{Populär} \\
\midrule
Emotion & Aktiviert Angst & Aktiviert Stolz \\
Vereinfachung & Lügt & Übersetzt \\
Feindbilder & Minderheiten & Strukturen \\
Resultat & Spaltet & Verbindet \\
\end{tabular}

\subsection*{Quick-Check vor jedem Auftritt}

\begin{itemize}
    \item[$\square$] Pivot vorbereitet? ($\rightarrow$ Kassenstellen)
    \item[$\square$] Flankierung für FPÖ-Swing? («nicht wer liest, sondern wer finanziert»)
    \item[$\square$] Steuergeld-Framing? (EUER Steuergeld, nicht MEHR)
    \item[$\square$] Passt es zu Babler? (Authentizitäts-Check)
\end{itemize}

\subsection*{Die goldene Regel}

\begin{center}
\fbox{\parbox{0.8\textwidth}{\centering\large\textbf{Nicht verteidigen. Pivotieren. Ordnen.}}}
\end{center}

\vfill

\begin{center}
{\color{fadarkblue}\rule{0.5\textwidth}{1pt}}

\vspace{0.5cm}

\textcopyright{} 2026 \textbf{\color{fadarkblue}FehrAdvice \& Partners AG}

\vspace{0.3cm}

{\small Behavioral Economics Consultancy Group}
\end{center}

\end{document}
